% sec5_1.tex — Section 5.1: The Error-Components Model

\section{The Error-Components Model}
\label{sec:5.1}

Our point of departure is the multiple-equation model with common coefficients
of Proposition~4.7. To simplify the discussion, we will assume that the sample is
a random sample. This assumption is satisfied in most popular panels, such as the
Panel Study of Income Dynamics (PSID) and the National Longitudinal Survey
(NLS). In those panels, a large number --- hundreds or even thousands --- of cross-section
units are randomly drawn from the population.

For this case of random samples, the assumptions of Proposition~4.7 can be
restated as (see (4.6.15) for the definition of $\mathbf{y}_i$ ($M \times 1$),
$\mathbf{Z}_i$ ($M \times L$), $\boldsymbol{\varepsilon}_i$ ($M \times 1$)):
\begin{align}
  &\text{(system of $M$ linear equations)}
  &\mathbf{y}_i &= \mathbf{Z}_i \boldsymbol{\delta} + \boldsymbol{\varepsilon}_i
  \label{eq:5.1.1} \tag{5.1.1} \\
  &\text{(random sample)}
  &\{\mathbf{y}_i, \mathbf{Z}_i\} &\text{ is i.i.d.}
  \label{eq:5.1.2} \tag{5.1.2} \\
  &\text{(``SUR assumption'')}
  &\E(\mathbf{z}_{im} \cdot \varepsilon_{ih}) &= \mathbf{0}
  \quad \text{for } m, h = 1, 2, \ldots, M,
  \label{eq:5.1.3} \tag{5.1.3}
\end{align}
i.e., $\E(\boldsymbol{\varepsilon}_i \otimes \mathbf{x}_i) = \mathbf{0}$ where
$\mathbf{x}_i = \text{union of } (\mathbf{z}_{i1}, \ldots, \mathbf{z}_{iM})$,
\begin{align}
  &\text{(identification)}
  &\E(\mathbf{Z}_i \otimes \mathbf{x}_i) &\text{ is of full column rank,}^{1}
  \label{eq:5.1.4} \tag{5.1.4} \\
  &\text{(conditional homoskedasticity)}
  &\E(\boldsymbol{\varepsilon}_i \boldsymbol{\varepsilon}_i' \mid \mathbf{x}_i)
  &= \E(\boldsymbol{\varepsilon}_i \boldsymbol{\varepsilon}_i') \equiv \boldsymbol{\Sigma}
  \label{eq:5.1.5} \tag{5.1.5} \\
  &\text{(nonsingularity of $\E(\mathbf{g}_i\mathbf{g}_i')$)}
  &\E(\mathbf{g}_i\mathbf{g}_i') &\text{ is nonsingular, where }
  \mathbf{g}_i \equiv \boldsymbol{\varepsilon}_i \otimes \mathbf{x}_i.
  \label{eq:5.1.6} \tag{5.1.6}
\end{align}

\footnotetext[1]{The $\boldsymbol{\Sigma}_{\mathbf{xz}}$ in Assumption~4.4$'$
can be written as $\E(\mathbf{Z}_i \otimes \mathbf{x}_i)$; see Review Question~4
to Section~4.6.}

As noted in Section~4.5, since
$\E(\mathbf{g}_i\mathbf{g}_i') = \boldsymbol{\Sigma} \otimes \E(\mathbf{x}_i\mathbf{x}_i')$
under conditional homoskedasticity, \eqref{eq:5.1.6} is equivalent to the condition
\begin{equation}
  \boldsymbol{\Sigma} \text{ and } \E(\mathbf{x}_i\mathbf{x}_i')
  \text{ are nonsingular.}
  \label{eq:5.1.6p} \tag{5.1.6$'$}
\end{equation}

Under these assumptions, the random-effects estimator $\hat{\boldsymbol{\delta}}_{\text{RE}}$,
defined in (4.6.8$'$) on page~293, is the efficient GMM estimator.

\subsection*{Error Components}

This model, very frequently used for panel data, is one where the unobservable
error term $\varepsilon_{im}$ is assumed to consist of two components:
\begin{equation}
  \varepsilon_{im} = \alpha_i + \eta_{im}.
  \label{eq:5.1.7} \tag{5.1.7}
\end{equation}

The first unobservable component $\alpha_i$, without the $m$ subscript and hence common
to all equations, is called the \textbf{individual effect}, the \textbf{individual heterogeneity}, or
the \textbf{fixed effect}. If we define
\[
  \underset{(M \times 1)}{\boldsymbol{\eta}_i}
  \equiv \begin{bmatrix} \eta_{i1} \\ \vdots \\ \eta_{iM} \end{bmatrix},
\]
the matrix representation of the $M$-equation system is
\[
  \begin{bmatrix} y_{i1} \\ \vdots \\ y_{iM} \end{bmatrix}
  = \begin{bmatrix} \mathbf{z}_{i1}'\boldsymbol{\delta} \\ \vdots \\ \mathbf{z}_{iM}'\boldsymbol{\delta} \end{bmatrix}
  + \begin{bmatrix} \alpha_i \\ \vdots \\ \alpha_i \end{bmatrix}
  + \begin{bmatrix} \eta_{i1} \\ \vdots \\ \eta_{iM} \end{bmatrix}
\]
or
\begin{equation}
  \mathbf{y}_i = \mathbf{Z}_i \boldsymbol{\delta} + \mathbf{1}_M \cdot \alpha_i
  + \boldsymbol{\eta}_i \quad (i = 1, 2, \ldots, n),
  \label{eq:5.1.1p} \tag{5.1.1$'$}
\end{equation}
where $\mathbf{1}_M$ is the $M$-dimensional vector of ones.

The orthogonality conditions \eqref{eq:5.1.3} are satisfied if the regressors of the system
are orthogonal to both error components, that is, if
\begin{align}
  \E(\mathbf{z}_{im} \cdot \alpha_i) &= \mathbf{0}
  \quad \text{for } m = 1, 2, \ldots, M,
  \label{eq:5.1.8a} \tag{5.1.8a} \\
  \E(\mathbf{z}_{im} \cdot \eta_{ih}) &= \mathbf{0}
  \quad \text{for } m, h = 1, 2, \ldots, M.
  \label{eq:5.1.8b} \tag{5.1.8b}
\end{align}

However, in many applications that utilize panel data, \eqref{eq:5.1.8a} is not a reasonable
assumption to make. This is because the fixed effect represents some permanent
characteristics of the individual economic unit, as illustrated by the following
examples.

\begin{example}[production function with firm heterogeneity]
\label{ex:5.1}
In the log-linear production function (3.2.13) of Section~3.2, suppose the firm's
efficiency $u_i$ stays constant over time. Then the equation for year $m$ is
\[
  \log(Q_{im}) = \phi_0 + \phi_1 \log(L_{im}) + u_i + v_{im},
\]
where $Q_{im}$ is the output of firm $i$ in year $m$, $L_{im}$ is labor input in year $m$, and
$v_{im}$ is the technology shock in year $m$. The equation can be written as \eqref{eq:5.1.1p}
by setting
\[
  y_{im} = \log(Q_{im}),\;
  \mathbf{z}_{im} = (1, \log(L_{im}))',\;
  \boldsymbol{\delta} = (\phi_0, \phi_1)',\;
  \alpha_i = u_i,\;
  \eta_{im} = v_{im}.
\]
Also,
\[
  \mathbf{x}_i = (1, \log(L_{i1}), \ldots, \log(L_{iM}))'.
\]

Under perfect competition, the individual effect $u_i$ would be positively correlated
with labor input because efficient firms, whose $u_i$ is higher than its
mean value, would hire more labor to expand (see (3.2.14)). If $v_{im}$ represents
unexpected shocks that are unforeseen by the firm when input choices are
made, it is reasonable to assume that $v_{im}$ is uncorrelated with the regressors.
\end{example}

\begin{example}[wage equation]
\label{ex:5.2}
For simplicity, drop experience from the
wage equation of Example~4.2 but suppose data for 1982, in addition to 1969
and 1980, are available. Suppose also that the coefficient of schooling and
that of IQ remain the same over time but that the intercept is time-dependent
(due, for example, to business-cycle effects on wages). With the error decomposition
\eqref{eq:5.1.7}, the three-equation system is
\[
  \begin{cases}
    LW69_i = \phi_1 + \beta S69_i + \gamma IQ_i + \alpha_i + \eta_{i1}, \\
    LW80_i = \phi_2 + \beta S80_i + \gamma IQ_i + \alpha_i + \eta_{i2}, \\
    LW82_i = \phi_3 + \beta S82_i + \gamma IQ_i + \alpha_i + \eta_{i3}.
  \end{cases}
\]

As noted at the end of Section~4.6, even if a subset of the coefficients differs
across equations, the system can be written as multiple equations with
common coefficients. In this example, this is accomplished with
\begin{equation}
  \mathbf{Z}_i = \begin{bmatrix}
    \mathbf{z}_{i1}' \\ \mathbf{z}_{i2}' \\ \mathbf{z}_{i3}'
  \end{bmatrix}
  = \begin{bmatrix}
    1 & 0 & 0 & S69_i & IQ_i \\
    0 & 1 & 0 & S80_i & IQ_i \\
    0 & 0 & 1 & S82_i & IQ_i
  \end{bmatrix}, \quad
  \boldsymbol{\delta}' = (\phi_1, \phi_2, \phi_3, \beta, \gamma).
  \label{eq:5.1.9} \tag{5.1.9}
\end{equation}
Also,
\[
  \mathbf{x}_i = (1, S69_i, S80_i, S82_i, IQ_i)'.
\]

The error term includes the wage determinants not accounted for by schooling
and IQ. It may be reasonable to assume that they are divided between
permanent individual traits (denoted $\alpha_i$), which would affect the individual's
choice of schooling, and other factors (denoted $\eta_{im}$) uncorrelated with the
regressors, such as measurement error in the log wage rate.
\end{example}

\subsection*{Group Means}

Fortunately, there is available a popular estimator, called the \textbf{fixed-effects estimator},
that is robust to the failure of \eqref{eq:5.1.8a}, i.e., that is consistent even when the
regressors are not orthogonal to the fixed effect $\alpha_i$. To anticipate the next section's
discussion, the estimator is applied to an $M$-equation system transformed from the
original system \eqref{eq:5.1.1p}. The matrix used for the transformation is the annihilator
associated with $\mathbf{1}_M$:
\begin{align}
  \underset{(M \times M)}{\mathbf{Q}}
  &\equiv \mathbf{I}_M - \mathbf{1}_M (\mathbf{1}_M' \mathbf{1}_M)^{-1} \mathbf{1}_M'
  \notag \\
  &= \mathbf{I}_M - \frac{1}{M} \mathbf{1}_M \mathbf{1}_M'
  \notag \\
  &= \mathbf{I}_M - \begin{bmatrix}
    1/M & \cdots & 1/M \\
    \vdots &        & \vdots \\
    1/M & \cdots & 1/M
  \end{bmatrix}.
  \label{eq:5.1.10} \tag{5.1.10}
\end{align}

What this matrix does is to extract \textbf{deviations from group means}. For example,
multiplying the $M$-dimensional vector $\mathbf{y}_i$ from left by $\mathbf{Q}$ results in
\begin{equation}
  \tilde{\mathbf{y}}_i \equiv \mathbf{Q}\mathbf{y}_i
  = \begin{bmatrix} y_{i1} - \bar{y}_i \\ \vdots \\ y_{iM} - \bar{y}_i \end{bmatrix}
  = \mathbf{y}_i - \mathbf{1}_M \cdot \bar{y}_i,
  \label{eq:5.1.11} \tag{5.1.11}
\end{equation}
where
\[
  \bar{y}_i = \frac{1}{M} \mathbf{1}_M' \mathbf{y}_i
  = \frac{1}{M} \sum_{m=1}^{M} y_{im}
\]
is the \textbf{group mean} for the dependent variable. Note that this mean is over $m$, \emph{not}
over $i$; it is specific to each group ($i$). Similarly for the regressors, each column of
$\mathbf{Q}\mathbf{Z}_i$ is the vector of deviations for the corresponding column of $\mathbf{Z}_i$.

\subsection*{A Reparameterization}

Not surprisingly, however, robustness to the failure of \eqref{eq:5.1.8a} comes at a price:
some of the parameters of the model may no longer be identifiable after the transformation
by $\mathbf{Q}$. To provide two examples,

\begin{itemize}
\item The obvious case is the coefficient of a variable common to all equations. For
  example, $IQ_i$ in Example~5.2 is a common regressor (see \eqref{eq:5.1.9}), so the column
  of $\mathbf{Z}_i$ corresponding to $IQ_i$ (the fifth column) is $\mathbf{1}_M \cdot IQ_i$ and the fifth
  column of $\mathbf{Q}\mathbf{Z}_i$ is a vector of zeros. Consequently, the IQ coefficient, $\gamma$, cannot
  be identified after transformation. To separate the coefficient of common
  regressors from the rest, let $\mathbf{b}_i$ be the vector of common regressors and write
  the $M \times L$ matrix of regressors, $\mathbf{Z}_i$, as
  \begin{equation}
    \mathbf{Z}_i = (\mathbf{F}_i \vdots \mathbf{1}_M \mathbf{b}_i')
    \label{eq:5.1.12} \tag{5.1.12}
  \end{equation}
  and partition the coefficient vector $\boldsymbol{\delta}$ accordingly:
  \[
    \boldsymbol{\delta} = \begin{bmatrix} \boldsymbol{\beta} \\ \boldsymbol{\gamma} \end{bmatrix}.
  \]
  The coefficient vector $\boldsymbol{\gamma}$ cannot be identified after transformation.

\item It may also be the case that, in addition to $\boldsymbol{\gamma}$, some of the coefficients in
  $\boldsymbol{\beta}$ are unidentifiable and a reparameterization is needed. Consider Example~5.2
  where each equation has a different intercept. One of the intercepts needs to be dropped
  from $\mathbf{F}_i$ and included in $\mathbf{b}_i$, because a common intercept can be created by taking
  a linear combination of the equation-specific intercepts. One reparameterization
  is to define
  \begin{equation}
    \mathbf{F}_i = \begin{bmatrix}
      1 & 0 & S69_i \\
      0 & 1 & S80_i \\
      0 & 0 & S82_i
    \end{bmatrix}, \quad
    \mathbf{b}_i = (1, IQ_i)',
    \label{eq:5.1.13} \tag{5.1.13}
  \end{equation}
  so
  \begin{equation}
    \mathbf{Z}_i = \begin{bmatrix}
      1 & 0 & S69_i & 1 & IQ_i \\
      0 & 1 & S80_i & 1 & IQ_i \\
      0 & 0 & S82_i & 1 & IQ_i
    \end{bmatrix}, \quad
    \boldsymbol{\beta} = (\phi_1 - \phi_3, \phi_2 - \phi_3, \beta)', \quad
    \boldsymbol{\gamma} = (\phi_3, \gamma)'.
    \label{eq:5.1.14} \tag{5.1.14}
  \end{equation}
\end{itemize}

More generally, the identification condition for fixed-effects estimation is that
$\mathbf{F}_i$ and $\mathbf{b}_i$ are defined so that
\begin{equation}
  \E(\mathbf{Q}\mathbf{F}_i \otimes \mathbf{x}_i) \text{ is of full column rank}
  \label{eq:5.1.15} \tag{5.1.15}
\end{equation}
where $\mathbf{x}_i = \text{union of } (\mathbf{z}_{i1}, \ldots, \mathbf{z}_{iM})$.
Why is this an identification condition?
Because (as you will see more clearly in Analytical Exercise~2), the fixed-effects
estimator is a specialization of (the) random-effects estimator applied to the transformed
regression of $\mathbf{Q}\mathbf{y}$ on $\mathbf{Q}\mathbf{F}$. \eqref{eq:5.1.15} is just the adaptation of the identification
condition \eqref{eq:5.1.4} to the transformed regression.

With $\mathbf{Z}_i$ thus divided between $\mathbf{F}_i$ and $\mathbf{b}_i$, the system of $M$-equations \eqref{eq:5.1.1p} on
page~325 can be rewritten as
\begin{align}
  \mathbf{y}_i &= \mathbf{F}_i \boldsymbol{\beta} + \mathbf{1}_M \cdot \mathbf{b}_i' \boldsymbol{\gamma}
  + \mathbf{1}_M \cdot \alpha_i + \boldsymbol{\eta}_i
  \quad (i = 1, 2, \ldots, n), \text{ or} \notag \\
  y_{im} &= \mathbf{f}_{im}' \boldsymbol{\beta} + \mathbf{b}_i' \boldsymbol{\gamma}
  + \alpha_i + \eta_{im}
  \quad (i = 1, \ldots, n;\; m = 1, \ldots, M),
  \label{eq:5.1.1pp} \tag{5.1.1$''$}
\end{align}
where $\mathbf{f}_{im}'$ is the $m$-th row of $\mathbf{F}_i$. The \textbf{error-components model} consists of assumptions
\eqref{eq:5.1.1pp} ($M$ linear equations), \eqref{eq:5.1.2} (random samples), \eqref{eq:5.1.8} (SUR
assumption with two error components), \eqref{eq:5.1.4} (identification), \eqref{eq:5.1.5} (conditional
homoskedasticity), \eqref{eq:5.1.6} (nonsingular $\E(\mathbf{g}_i\mathbf{g}_i')$), and \eqref{eq:5.1.15} (fixed-effects
identification). We include the additional identification condition \eqref{eq:5.1.15} for fixed-effects
estimation, so that both the random-effects estimator and the fixed-effects
estimator are well-defined for the same model.

\bigskip
\noindent\rule{\textwidth}{0.5pt}
\textsc{Questions for Review}
\medskip

\begin{enumerate}
\item (Verification of Proposition~4.7 assumptions)\quad
  Verify that the conditions of
  Proposition~4.7 are satisfied when \eqref{eq:5.1.1}--\eqref{eq:5.1.6} hold. Show that the identification
  condition \eqref{eq:5.1.4} is satisfied when $\E(\mathbf{x}_i\mathbf{x}_i')$ is nonsingular.
  \textbf{Hint:} $\mathbf{z}_{im}$ is a subset of $\mathbf{x}_i$. Thus, \eqref{eq:5.1.4} is redundant.

\item (Nonuniqueness of reparameterization)\quad
  \eqref{eq:5.1.13} is not the only reparameterization
  for Example~5.2. If $\mathbf{F}_i$ is defined as
  \[
    \mathbf{F}_i = \begin{bmatrix}
      0 & 0 & S69_i \\
      1 & 0 & S80_i \\
      0 & 1 & S82_i
    \end{bmatrix},
  \]
  how should $\mathbf{b}_i$ and $\boldsymbol{\delta}$ be defined?

\item Without reparameterization, the $\mathbf{F}_i$ for Example~5.2 is
  \[
    \mathbf{F}_i = \begin{bmatrix}
      1 & 0 & 0 & S69_i \\
      0 & 1 & 0 & S80_i \\
      0 & 0 & 1 & S82_i
    \end{bmatrix}
  \]
  with $\mathbf{b}_i = IQ_i$. Verify that \eqref{eq:5.1.15} is not satisfied.
\end{enumerate}
